% !TEX root = ./main.tex

\section{Related Work}\label{sec:related_work}
 
 
\paragraph{Expectation Transformers}
Nondeterminism in probabilistic programs with loops and unbounded numbers (but without conditioning) has been studied by \citet{mciver2005abstraction}, through \textit{expectation transformers}.
This work was later extended by \citet{olmedo2018conditioning} to discrete probabilistic programs with conditioning and loops but without nondeterminism.
Additionally,  \citet{olmedo2018conditioning}  showed that nondeterminism and conditioning cannot both be covered by \emph{inductive} expectation transformers.
They (and previously \citet{shoup2009computational}) also showed that loops can model observations, essentially by emulating rejection sampling.
However, this requires reasoning about loops, which is particularly difficult in the presence of randomness~\cite{DBLP:journals/acta/KaminskiKM19}.

\paragraph{Probabilistic Programming Languages \textnormal{(PPLs)}}
A variety of PPLs for Bayesian reasoning has been developed.
In contrast to our work, these often focus on continuous distributions.
In fact, some PPLs like Aqua~\cite{huang2021aqua}, Pyro's Variational Inference~\cite{bingham2019pyro}, or Stan~\cite{carpenter2017stan} \emph{only} support continuous distributions.
Nonetheless, there are many languages which support discrete random variables alongside continuous distributions, for example Anglican~\cite{wood-aistats-2014}, Hakaru~\cite{narayanan2016probabilistic}, SOGA~\cite{randone2024inference} or WebPPL~\cite{dippl}.
However, \emph{none of the above PPLs support nondeterminism}.

A notable PPL which allows nondeterminism via decision making is ProbLog~\cite{fierens2015inference}. 
ProbLog features both conditioning via evidence and an extension to determine optimal decisions~\cite{van2010dtproblog}, but explicitly forbids combining them.
Another example is IBAL~\cite{pfeffer2007design}, which includes a choose-statement to pick the \textit{locally} optimal value to maximize a utility function.
However, \Cref{sec:subOpt} shows that local reasoning can be insufficient. 

Languages like Psi~\cite{gehr2020lambdapsi} allow unspecified inputs, resulting in parameterized inference.
Parameters are fundamentally different from nondeterminism.
In particular, parameters are fixed at program start and hence the semantics of \ref{ExLet1} and \ref{ExLet2} cannot easily be distinguished.

PMAF~\cite{wang2018pmaf} is a framework for the analysis of probabilistic programs with unstructured control flow, recursion, nondeterminism, and observations.
PMAF can be instantiated with a variety of analyses, and it may be possible to perform \NDICE inference within the framework. 
However, the semantics of observations differ between PMAF and \NDICE \cite{wang2018denotational}:
Both assign probability $0$ to failed observations, but PMAF does not normalize the resulting distributions.
Hence, the PMAF equivalent to $\OBSERVE{\FLIP{\frac{2}{3}}}$ only returns $\TRUE$ with probability $\frac{2}{3}$.

\paragraph{Probabilistic Model Checking and Bisimilarity}
Since \NDICE reduces inference to a conditional reachability property on an MDP, every problem which can be modeled by a \NDICE program could also be modeled directly as an MDP.
Then, a \textit{probabilistic model checker} can be used to perform model checking as was seen in the empirical evaluation.
This approach offers many advantages compared to \NDICE:
For one, a variety of mature model checkers has been developed over the years \cite{hahn20192019}.
Secondly, MDPs themselves are more expressive than \NDICE programs as they can model cycles, concurrency and symbolic unbounded numbers.
Leveraging these advantages, \citet{jansen2016bounded} performed inference on probabilistic programs with nondeterminism by simply unrolling the operational semantics of programs as an MDP, and letting a probabilistic model checker determine the conditional reachability probabilities.
\Cref{sec:implementation} shows that this can be an effective approach in practice, but the constructed MDPs tend to be very large.

Hence, it is reasonable to consider approaches like \NDICE, which aim to reduce the size of the final MDP in order to speed up model checking.
Here, model checkers also use ADDs \cite{baier1997symbolic}, but contrary to our work ADDs are used to compactly represent the full state space, rather than to remove states.
Still, probabilistic model checkers have developed a variety of other techniques to speed up model checking \cite{hartmanns2025revised}.
Here, \emph{probabilistic bisimulation} \cite{larsen1989bisimulation, segala1995probabilistic} is particularly relevant to our work, since it presents a generic way to minimize MDPs.
Most notably, the redundant states removed by \NDICE and bisimulations differ, and both approaches remove redundancies which the other misses:
Bisimulations remove symmetries such as e.g. distinct flip operations which use the same probabilities, whereas the decision diagrams of \NDICE remove \enquote{deterministic} states with only a single successor.

Regarding the empirical evaluation in \Cref{sec:implementation}, we found that using bisimulations to reduce the directly encoded MDPs used by Storm did not yield any speedups compared to applying Storm directly.
In fact, calculating the minimized MDP was usually very expensive, and frequently resulted in time-outs.
Consequently, we used Storm without minimizing the MDP first.


\paragraph{MDP Modeling Languages}
	In some ways, \NDICE can be seen as a modeling language for (acyclic) MDPs, and so we will compare it to existing MDP modeling languages like Jani~\cite{budde2017jani} or Prism~\cite{KNP11}.
	Jani and Prism are used to specify the input for model checking tools like Storm \cite{hensel2022probabilistic} or Modest \cite{hartmanns2014modest}, and they allow one to define general MDPs which may contain cycles or be constructed from multiple modules running in parallel. 
	Regarding modeling itself, \NDICE specifies models with the syntax and structure of programs, Prism uses a state-based language whereas Jani encodes models in the JSON data format. 
	However, both Prism and Jani use the same modeling approach:	
	
	First, a set of variables is specified, and the state space of the MDP is defined as the set of possible variable assignments. 
	Transitions then manipulate the variables, and their effects are defined through symbolic expressions which specify the updated values of the variables.
	Whether or not a transition is enabled in a state is defined with a guard, i.e. a boolean condition on the variables. 
	A transition can only be taken when its guard condition holds for the state, but beyond that transitions are always enabled.
	This is useful when modeling e.g. a server which should always be waiting for different requests, but is inconvenient for scenarios where events have to occur in a strict order, such as \Cref{fig:plane}.
	 Hence, MDPs often need an additional state variable to explicitly capture which transitions should be enabled next, see for example \citet{d1997bounded}.
	 In contrast, \NDICE automatically defines the order of operations through the natural program order.



\paragraph{Marginal MAP}
The goal of a marginal maximum a posteriori (MMAP) query is closely related to our inference task: 
Find an optimal assignment to some variables -- which can be interpreted as a nondeterministic choice -- in order to maximize the marginal probability over the remaining variables \cite{liu2013variational, choi2022solving}.
Contrary to this very similar problem statement, we believe that MMAP problems and \NDICE inference require distinct solution approaches.

For one, \NDICE inference cannot easily be performed by MMAP solvers, as the goal of MMAP queries is to find an optimal \textit{assignment} to variables.
An assignment fixes a single value for every nondeterministic choice, which in the context of \NDICE programs means statically replacing every $\NFLIP$ statement with $\TRUE$ or $\FALSE$.
Our model of nondeterminism is more general, hence formulating \NDICE inference as a MMAP query may require enumerating all contexts for $\NFLIP$ as a separate choice, leading to exponentially larger input programs.
%
Similarly, \NDICE cannot easily solve MMAP queries, as \NDICE may consider too many behaviors.
Still, MMAP queries could be compiled into \NDICE programs by performing all nondeterministic choices at the start of the program to avoid acting based on probabilistic outcomes. 
However, this removes some of the program structure which is exploited by ADDs, which would negatively affect the inference efficiency.



\section{Conclusion and Future Work}\label{sec:future}
We have introduced syntax and semantics of \NDICE, a probabilistic programming language which conservatively extends \DICE with nondeterminism.
We also presented a novel inference algorithm for probabilistic programs with nondeterminism.
The algorithm compiles programs to Boolean formulas and reduces inference queries to conditional reachability in an MDP.
We evaluated a prototype implementation and showed that our algorithm can feasibly do inference for large programs.
We believe our work is a promising step towards efficient inference for probabilistic programs with nondeterminism and conditioning.
As future work, we deem the following worthwhile.

\paragraph{Unbounded Loops}
PPLs like \DICE, \NDICE or Psi \cite{gehr2016psi} can only approximate the behavior of unbounded loops with bounded loops.
However, \citet{torres2024iteration} extended \DICE with a statement for \textit{unbounded} iteration, as well as an exact inference method for programs with this new statement.
We are eager to see if their work can be extended with nondeterminism.

\paragraph{Determining Optimal Programs}
Our work only determines the maximum probability to return a value, but it may also be of interest to retrieve a purely probabilistic program which generates those probabilities.
Synthesizing optimal deterministic programs has been considered without conditioning in \citet{batz2024programmatic}, but it is unclear if this approach can be generalized.

